% Independent mathematical audit passed: random_hosts. Dependencies separately audited.
\section{The sharp order for cyclic cuts of size at most four}
\label{sec:four-cut-sharp-order}

The linear two-cycle packing bound improves both four-cut estimates.
For aligned saturated shores, the largest forest terms cancel in the gluing identity.
These two improvements give the same order as the cyclic-square lower bound.

\begin{corollary}[A linear four-pole aggregate error]
\label{cor:linear-four-pole-aggregate}
There is an absolute constant $A'_4$ such that every connected loopless cubic four-pole of rank $h$ satisfies
\[
 \sum_{i<j}P_{ij}+\sum_\alpha f_\alpha
 \le4\,2^h+A'_4(h+3)2^{h/2}.
\]
\end{corollary}

\begin{proof}
Repeat the proof of Theorem~\ref{thm:four-pole-aggregate}.
Its only fourth-degree factor comes from the cross-sum forest estimate.
Replace that estimate by Corollary~\ref{cor:linear-four-pole-cross-sum-forest}.
The low-cross-sum full-star case still has error $O(2^{h/2})$.
All fixed rank thresholds and the finitely many smaller ranks are absorbed into $A'_4$.
\end{proof}

\begin{theorem}[Sharp-order four-cut error]
\label{thm:sharp-order-four-cut-error}
There is an absolute constant $K$ such that every connected loopless cubic multigraph of rank $r$ with a cyclic edge cut of size at most four satisfies
\[
 c(G)\le2^{r-1}+Kr2^{r/2}.
\]
The factor $r$ cannot be replaced by a function that is $o(r)$ throughout this class.
\end{theorem}

\begin{proof}
Cuts of size at most three already have error $O(2^{r/2})$ by Theorem~\ref{thm:all-three-cut-half-exponent}.
For a cyclic four-edge cut, use the connected four-pole shores from Theorem~\ref{thm:all-four-cut-half-exponent}.
Write their ranks as $h,k$, so $r=h+k+3$.

First suppose both ranks exceed the fixed threshold in Theorem~\ref{thm:four-cut-large-shores}.
Increase that threshold if needed for the constant in the linear forest estimate.
Its ordinary and nonaligned exceptional cases have fixed crossing deficits.
Those deficits still absorb all internal cycles.
The same profile classification applies when the forest estimate is replaced by its linear refinement.
Only equal full-star supports and equal full-$K_{2,2}$ supports need a smaller error bound.

For equal stars, the earlier bound is
\[
 \mathcal B\le4+3\rho_h\rho_k,
 \qquad \rho_j=C2^{-j/2}.
\]
Thus their crossing error is $O(2^{(h+k)/2})$.
Both shores have a full terminal pair, so their internal cycle counts have the same half-exponent bound.

Now suppose the two full-$K_{2,2}$ supports agree.
They have the same distinguished same-side cap pairing $\alpha$.
For the rank-$h$ shore, put
\[
 C_h=f_\alpha,\qquad D_h=\sum_{\beta\ne\alpha}f_\beta.
\]
Thus $D_h$ counts through pairs and $C_h$ counts pairs of disjoint same-side caps.
Define $C_k,D_k$ in the other shore.
The audited saturated-block bounds give
\[
 D_h\le2^{(h+1)/2},\qquad
 C_h\le3(h+1)2^{(h+1)/2},
\]
and the corresponding inequalities at rank $k$.

The exact four-cut gluing identity subtracts products with equal pairings.
In particular, its $C_hC_k$ term cancels.
The remaining four-edge contribution is at most
\[
 C_hD_k+D_hC_k+D_hD_k
 \le[3(h+k+2)+1]2^{(h+k+2)/2}.
\]
The two nonfull path-coordinate products contribute only $O(2^{(h+k)/2})$.
Internal cycles again contribute $O(2^{h/2}+2^{k/2})$.
This proves the required $O(r2^{r/2})$ bound when both shores have large rank.

Suppose instead that one shore has bounded rank.
The proof of Theorem~\ref{thm:all-four-cut-half-exponent} already places every low-path-density large shore below the one-half bound.
Its only remaining case uses the universal four-pole aggregate and a high path count in the large shore.
Replace the aggregate there by Corollary~\ref{cor:linear-four-pole-aggregate}.
The resulting crossing error is $O((k+3)2^{k/2})$, since the other rank is bounded.
The large shore's internal cycles are $O(2^{k/2})$ by the fixed-density tradeoff.
All bounded ranks are absorbed into $K$.
This proves the upper bound.

For the lower bound, take the cyclic-square graph $G_q$ from Proposition~\ref{prop:hk-short-path-obstruction}.
It has rank $r=2q+1$ and a cyclic four-edge cut around each square.
There are $4^q$ cycles using one edge at every interface.
There are $q$ cycles using no interface.
A proper interval of $d$ interfaces gives $q2^{d+1}$ cycles, for $1\le d\le q-1$.
Using both edges at every interface gives $2^{q-1}$ cycles.
Therefore
\[
 \begin{split}
 c(G_q)
 &=4^q+q+q\sum_{d=1}^{q-1}2^{d+1}+2^{q-1}\\
 &=2^{2q}+\left(2q+\frac12\right)2^q-3q.
 \end{split}
\]
Its excess is $\Theta(q2^q)=\Theta(r2^{r/2})$.
\end{proof}

This theorem addresses a structural subclass of cubic kernels.
It does not prove the unrestricted sharp-order error.
After the established reductions, that remaining question concerns cyclically five-edge-connected cubic kernels.

The independent script \texttt{check\_four\_cut\_sharp\_order.py} verifies the total cycle formula for $2\le q\le7$ by binary cycle-space enumeration.
It also verifies the disjoint-cap packing formula for one through six square blocks.
Its results are saved in \texttt{four-cut-sharp-order-checks.json}.
